\uuid{dypq}
\titre{Matrice/Transposée/Inversibilité/Rang}
\auteur{hanani ; mzari}
\organisation{exo7}
\licence{CC BY-NC-SA 4.0}
\datecreate{2020-01-20}
\niveau{L1}
\module{Algèbre}
\chapitre{Matrice}
\sousChapitre{Inverse, méthode de Gauss}
\difficulte{3}
\theme{Inverse d'une matrice, Matrice, Transposée, Inversibilité, Rang}
\exo7qcmid{lille-2.matrices.029}

\contenu{
\begin{question}
\qclassification{108.05}
\qtags{motcle=Matrice/Transposée/Inversibilité/Rang}
Soit $A$ et $B$ deux matrices carrées d'ordre $n\ge 1$. On notera $^tA$ la transposée de $A$ et $\mbox{rg}\, (A)$ le rang de $A$.
Quelles sont les assertions vraies ?
\begin{answers}
\good{$\mbox{rg}(A)=\mbox{rg}( \, ^tA)$}.
\good{Si $A$ est inversible, $\mbox{rg}(A)=\mbox{rg}(A^{-1})$}.
\bad{$\mbox{rg}(A+B)=\max \big(\mbox{rg}(A), \mbox{rg}(B)\big)$}.
\bad{$\mbox{rg}(AB)= \mbox{rg}(BA)$}.
\end{answers}
\begin{explanations} Le rang d'une matrice est le nombre maximum de vecteurs colonnes ou lignes qui sont linéairement indépendants. On en déduit que $\mbox{rg}(A)=\mbox{rg}(\, ^tA)$ et que, si $A$ est inversible, $\mbox{rg}(A)=\mbox{rg}(A^{-1})=n$.
\vskip0mm
En général, $\mbox{rg}(A+B)\neq \max \big(\mbox{rg}(A), \mbox{rg}(B)\big)$ et $\mbox{rg}(AB)\neq \mbox{rg}(BA)$. Contre-exemple : avec
$$A = \left(\begin{array}{rc}
1&-1\\-1&1\\ \end{array}\right),\quad B = \left(\begin{array}{rc}
-1&1\\ 1&-1\\ \end{array}\right),\quad C =
\left(\begin{array}{rc} 1&2\\1&2\\ \end{array}\right),$$
on vérifie que :  $A+B =0$, $AC =0$ et $CA=
\left(\begin{array}{rc}-1&1\\
-1&1\\ \end{array}\right)$. Donc $\mbox{rg}(A)=\mbox{rg}(B)=1$ mais $\mbox{rg}(A+B)=0$ et $\mbox{rg}(AC)= 0$ est différent de $\mbox{rg}(CA)=1$.
\end{explanations}
\end{question}
}
